\section{Test}
\subsection{NVDA}
L'uso di NVDA ci ha permesso di identificare problemi che altrimenti non avremmo notato, qui di seguito vengono riportati i problemi riscontrati e le soluzioni adottate:
\begin{itemize}
    \item Le dialog si chiudevano automaticamente quando il tab raggiungeva un elemento al di fuori della dialog.\\
    \textit{Soluzione:} abbiamo utilizzato javascript per limitare la chiusura della dialog solo al click della x.
    \item Nelle tabelle il contenuto descrittivo, associato alla tabella tramite l'attributo aria-describedby, veniva letto due volte.\\
    \textit{Soluzione:} abbiamo aggiunto l'attributo aria-hidden al tag che contiene la descrizione, in questo modo il contenuto descrittivo viene letto solo una volta incontrato l'attributo aria-describedby.
\end{itemize}

\subsection{Total Validator}
Total Validator ci ha permesso di identificare errori di sintassi HTML e errori di accessibilità. Abbiamo svolto i test due volte: la prima volta durante lo sviluppo, per scoprire errori prima che venissero propagati, e la seconda volta a fine sviluppo, per assicurarci che non ci fossero errori residui.\\
\subsection{Test vari con Silktide}
L'uso dell'estensione Silktide, reperibile gratiutamente su Chrome al seguente link \url{https://chromewebstore.google.com/detail/mpobacholfblmnpnfbiomjkecoojakah?utm_source=item-share-cb}, è stata utilizzata solo per simulare l'esperienza di utenti con specifiche disabilità visive, come il datonismo e la cataratta oppure per simulare utenti con disabilità cognitive come la dislessia.\\
Inoltre Silktide si è dimostrato utile anche per rendere evidenti caratteristiche strutturali come la gerarchia dei titoli e l'ordine del focus, che altrimenti avremmo dovuto verificare manualmente e con forse scarsa efficacia.\\

\subsection{Test con Lighthouse}
Lighthouse è uno strumento di Google che ci ha permesso di valutare le prestazioni. Il suo uso è stato fondamentale per accorgerci del problema del peso delle immagini, che rallentava notevolmente il caricamento del sito.\\
Dopo i test con lighthouse, di cui ora riportiamo gli ottimi risultati di una pagina critica, abbiamo deciso di ridurre il peso delle immagini cambiando la funzione di upload per ridurne il peso di quelle che vengono caricate dagli utenti (amministratori o non) a meno di 40KB.\\
\begin{figure}[H]
    \centering
    \includegraphics[width=0.4\textwidth]{assets/lighthouse-animali.png}
    \caption{Risultati Lighthouse per la pagina animali}
\end{figure}
\subsection{Altri test}
Qui di seguito riportiamo una lista di altri test che abbiamo svolto durante lo sviluppo del sito web:
\begin{itemize}
    \item test su più browser;
    \item test reponsività su più dispositivi (desktop, tablet/ipad, smartphone);
    \item test sull'utenza, con persone di diverse età e con diverse competenze digitali;
    \item test di inserimento dati non validi o che non rispettano i vincoli del database, per verificare la robustezza delle operazioni sul database.
\end{itemize}